\documentclass[11pt]{book}

\usepackage[utf8]{inputenc}
\usepackage[T1]{fontenc}
\usepackage{lmodern}
\usepackage[margin=2.5cm]{geometry}
\usepackage{hyperref}
\usepackage{listings}
\usepackage{xcolor}
\usepackage{enumitem}
\usepackage{titlesec}

% --- Kod listingi stili ---------------------------------------------------
\definecolor{codebg}{RGB}{245,245,245}
\definecolor{codekey}{RGB}{0,102,153}
\definecolor{codestr}{RGB}{163,21,21}
\definecolor{codecom}{RGB}{0,128,0}

\lstdefinestyle{base}{
  backgroundcolor=\color{codebg},
  basicstyle=\ttfamily\small,
  keywordstyle=\color{codekey}\bfseries,
  stringstyle=\color{codestr},
  commentstyle=\color{codecom}\itshape,
  breaklines=true,
  showstringspaces=false,
  frame=single,
  rulecolor=\color{black!20},
  numbers=left,
  numberstyle=\tiny\color{black!50},
  numbersep=8pt,
  tabsize=2,
}
\lstdefinelanguage{yaml}{
  keywords={services,image,environment,ports,depends_on,volumes,build},
  sensitive=true,
  comment=[l]{\#},
  string=[s]{'}{'},
}

\hypersetup{
  colorlinks=true,
  linkcolor=blue!60!black,
  urlcolor=blue!60!black,
  citecolor=blue!60!black,
}

\title{REST API and Modern Web Services\\
\large A Practical Companion to the \texttt{rest-api-book} Repository}
\author{İbrahim Güney}
\date{\today}

\begin{document}

\frontmatter
\maketitle

\tableofcontents

\chapter*{Preface}
\addcontentsline{toc}{chapter}{Preface}

This book is the written companion to a live GitHub repository
(\texttt{rest-api-book}) that implements a small but realistic web service:
a FastAPI application backed by PostgreSQL, protected by JWT-based
authentication, packaged with Docker, and validated by an automated test
suite. Every chapter ends with pointers to the corresponding files in the
repository so that what is described in prose can be inspected, executed,
and modified.

\mainmatter

% =========================================================================
\chapter{Introduction}
\label{ch:intro}

Modern applications rarely live in isolation. A mobile client, a web
front-end, and a third-party integration may all need to talk to the same
core data. The most widely deployed pattern for that conversation is the
\emph{REST API}, an architectural style that uses HTTP verbs, resource
URIs, and stateless requests to expose application functionality.

\section{Goals of the Book}
By the end of this book, the reader should be able to:
\begin{itemize}[leftmargin=*]
  \item explain the constraints of REST and apply them to API design;
  \item build a typed Python web service with FastAPI;
  \item persist data in PostgreSQL through SQLAlchemy;
  \item implement password hashing and JWT-based authentication;
  \item containerize the service with Docker and Docker Compose;
  \item write meaningful tests with \texttt{pytest} and FastAPI's
        \texttt{TestClient};
  \item maintain the project on GitHub with a sustainable workflow.
\end{itemize}

\section{Repository Map}
The companion repository is organized as follows:
\begin{lstlisting}[style=base]
rest-api-book/
├── api/
│   ├── auth.py          # JWT, password hashing, current-user dependency
│   ├── database.py      # SQLAlchemy engine + session
│   ├── main.py          # FastAPI app entry point
│   ├── models.py        # ORM models
│   ├── schemas.py       # Pydantic request/response schemas
│   └── routes/
│       ├── students.py  # /students CRUD
│       └── users.py     # /users register, login, me
├── tests/               # pytest suite
├── book/main.tex        # this book
├── Dockerfile           # API image
├── docker-compose.yml   # API + PostgreSQL
└── requirements.txt
\end{lstlisting}

% =========================================================================
\chapter{REST Fundamentals}
\label{ch:rest}

REST (\emph{Representational State Transfer}) was introduced by Roy
Fielding in his 2000 doctoral dissertation. It is not a protocol but a set
of architectural constraints that, when satisfied, give a distributed
system the properties of scalability, evolvability, and visibility.

\section{The Six Constraints}
\begin{enumerate}[leftmargin=*]
  \item \textbf{Client--server.} Clear separation of concerns between the
        user interface and data storage.
  \item \textbf{Stateless.} Each request carries everything the server
        needs to process it; the server keeps no client session.
  \item \textbf{Cacheable.} Responses must declare whether they are
        cacheable, so that intermediaries can short-circuit traffic.
  \item \textbf{Uniform interface.} Resources are identified by URIs,
        manipulated through representations, and described by
        self-descriptive messages.
  \item \textbf{Layered system.} A client cannot tell whether it is
        talking to the origin server or an intermediary.
  \item \textbf{Code on demand} (optional). The server may extend client
        functionality by transferring executable code.
\end{enumerate}

\section{HTTP Verbs as Resource Operations}
\begin{tabular}{l l l}
\textbf{Verb} & \textbf{Meaning} & \textbf{Idempotent?} \\\hline
GET    & Retrieve a representation        & Yes \\
POST   & Create a new resource            & No  \\
PUT    & Replace a resource               & Yes \\
PATCH  & Partially update a resource      & No  \\
DELETE & Remove a resource                & Yes \\
\end{tabular}

\section{Status Codes That Matter}
A REST API communicates outcomes with status codes. The minimum viable
vocabulary is:
\textbf{200 OK}, \textbf{201 Created}, \textbf{204 No Content},
\textbf{400 Bad Request}, \textbf{401 Unauthorized},
\textbf{403 Forbidden}, \textbf{404 Not Found},
\textbf{409 Conflict}, \textbf{422 Unprocessable Entity},
\textbf{500 Internal Server Error}.

\section{Designing Resources}
A good URL describes a \emph{noun}, not a verb:
\texttt{/students/42} is preferable to \texttt{/getStudentById?id=42}.
Nested resources express containment
(\texttt{/students/42/courses}), while query parameters express
filters and pagination (\texttt{/students?department=cs\&limit=20}).

% =========================================================================
\chapter{FastAPI in Practice}
\label{ch:fastapi}

FastAPI is a modern Python framework that builds on Starlette for the web
layer and Pydantic for data validation. Its key promise is that the same
type hints that make Python code self-documenting also generate request
validation, response serialization, and an OpenAPI schema.

\section{The Application Object}
The repository's entry point is \texttt{api/main.py}:
\begin{lstlisting}[style=base, language=Python]
from fastapi import FastAPI
from api.database import Base, engine
from api import models  # noqa: F401
from api.routes import students, users

Base.metadata.create_all(bind=engine)

app = FastAPI(title="REST API Book")
app.include_router(students.router, prefix="/students", tags=["Students"])
app.include_router(users.router,    prefix="/users",    tags=["Users"])
\end{lstlisting}

\section{Routers, Dependencies, and Tags}
\texttt{APIRouter} groups related endpoints. The \texttt{Depends()} marker
expresses the dependencies of an endpoint --- a database session, the
current user, a query parameter --- and FastAPI resolves them on every
request. Tags drive the grouping that the auto-generated documentation
displays at \texttt{/docs}.

\section{Pydantic Schemas}
Request and response shapes live in \texttt{api/schemas.py}. By providing
both an \texttt{StudentCreate} (input) and a \texttt{StudentRead} (output)
class we keep the persistence layer decoupled from what the API exposes.
The \texttt{model\_config = ConfigDict(from\_attributes=True)} line tells
Pydantic that it can build a response straight from a SQLAlchemy ORM
object.

\section{Automatic OpenAPI}
Once the application is running, the URL \texttt{/openapi.json} returns a
machine-readable description of every endpoint, and \texttt{/docs} renders
the interactive Swagger UI. Clients in any language can be generated from
that schema.

% =========================================================================
\chapter{PostgreSQL and SQLAlchemy}
\label{ch:db}

PostgreSQL is the project's database of record. SQLAlchemy provides two
layers: a low-level \emph{Core} (SQL expression language) and an
\emph{ORM} that maps Python classes to tables. The repository uses the
ORM.

\section{Engine and Session}
The engine is a connection pool with dialect awareness. Sessions are
short-lived units of work obtained from a session factory:
\begin{lstlisting}[style=base, language=Python]
DATABASE_URL = os.getenv(
    "DATABASE_URL",
    "postgresql://admin:password@db:5432/university",
)
engine = create_engine(DATABASE_URL)
SessionLocal = sessionmaker(bind=engine, autocommit=False, autoflush=False)
\end{lstlisting}

\section{Declarative Models}
Each model inherits from \texttt{Base} and declares its columns:
\begin{lstlisting}[style=base, language=Python]
class Student(Base):
    __tablename__ = "students"
    id         = Column(Integer, primary_key=True, index=True)
    name       = Column(String, nullable=False)
    department = Column(String, nullable=False)
\end{lstlisting}

\section{The \texttt{get\_db} Dependency}
Endpoints obtain a session through a generator-based dependency:
\begin{lstlisting}[style=base, language=Python]
def get_db():
    db = SessionLocal()
    try:
        yield db
    finally:
        db.close()
\end{lstlisting}
This pattern guarantees the session is closed even if the endpoint
raises.

\section{Migrations}
For a teaching project we rely on
\texttt{Base.metadata.create\_all()} at startup. Real systems should
manage schema changes with a tool such as Alembic.

% =========================================================================
\chapter{Authentication with JWT}
\label{ch:auth}

Authentication answers \emph{who is calling}; authorization answers
\emph{are they allowed}. The repository implements stateless
authentication with JSON Web Tokens.

\section{Hashing Passwords}
Plain passwords are never stored. \texttt{passlib} with the
\texttt{bcrypt} backend computes a salted hash:
\begin{lstlisting}[style=base, language=Python]
pwd_context = CryptContext(schemes=["bcrypt"], deprecated="auto")
def hash_password(p):    return pwd_context.hash(p)
def verify_password(p,h): return pwd_context.verify(p, h)
\end{lstlisting}

\section{Issuing a Token}
On a successful login, the server signs a JWT that carries the user's
email as the \texttt{sub} claim:
\begin{lstlisting}[style=base, language=Python]
expire = datetime.now(timezone.utc) + timedelta(minutes=30)
payload = {"sub": user.email, "exp": expire}
token = jwt.encode(payload, SECRET_KEY, algorithm="HS256")
\end{lstlisting}

\section{Verifying on Each Request}
Protected endpoints depend on \texttt{get\_current\_user}, which decodes
the bearer token and loads the matching user from the database:
\begin{lstlisting}[style=base, language=Python]
def get_current_user(token: str = Depends(oauth2_scheme),
                     db: Session = Depends(get_db)) -> User:
    try:
        payload = jwt.decode(token, SECRET_KEY, algorithms=[ALGORITHM])
        email = payload.get("sub")
    except JWTError:
        raise HTTPException(401, "Could not validate credentials")
    user = db.query(User).filter(User.email == email).first()
    if not user:
        raise HTTPException(401, "Could not validate credentials")
    return user
\end{lstlisting}
Any route that adds \texttt{\_: User = Depends(get\_current\_user)} to its
signature is automatically protected.

\section{Operational Notes}
The signing key is read from \texttt{SECRET\_KEY}. In development the
default value is fine; in production it \emph{must} be supplied through
the environment and rotated on suspicion of compromise.

% =========================================================================
\chapter{Containerization with Docker}
\label{ch:docker}

\section{The API Image}
\texttt{Dockerfile} builds a single-purpose image based on the official
\texttt{python:3.11} runtime:
\begin{lstlisting}[style=base]
FROM python:3.11
WORKDIR /app
COPY requirements.txt .
RUN pip install -r requirements.txt
COPY . .
CMD ["uvicorn", "api.main:app", "--host", "0.0.0.0", "--port", "8000"]
\end{lstlisting}

\section{Compose Topology}
\texttt{docker-compose.yml} wires the API to a PostgreSQL container and
declares a named volume so that data survives a restart:
\begin{lstlisting}[style=base, language=yaml]
services:
  api:
    build: .
    ports: ["8000:8000"]
    environment:
      DATABASE_URL: postgresql://admin:password@db:5432/university
      SECRET_KEY:   ${SECRET_KEY:-change-this-secret-key}
    depends_on: [db]
  db:
    image: postgres:15
    environment:
      POSTGRES_USER:     admin
      POSTGRES_PASSWORD: password
      POSTGRES_DB:       university
    ports: ["5432:5432"]
    volumes: ["postgres_data:/var/lib/postgresql/data"]
volumes:
  postgres_data:
\end{lstlisting}

\section{Running the Stack}
A single command brings everything up:
\begin{lstlisting}[style=base]
docker compose up --build
\end{lstlisting}
The API is then available at \url{http://localhost:8000} and the
documentation at \url{http://localhost:8000/docs}.

% =========================================================================
\chapter{Testing}
\label{ch:tests}

A REST API without tests is a contract whose terms can change silently.
The repository uses \texttt{pytest} and FastAPI's \texttt{TestClient},
which speaks to the application in-process without needing a real HTTP
server.

\section{Isolating the Test Database}
\texttt{tests/conftest.py} sets \texttt{DATABASE\_URL} to a local SQLite
file \emph{before} the application is imported. SQLAlchemy then builds a
SQLite engine, and tables are recreated for each test. The production
PostgreSQL configuration is untouched.

\section{Fixtures}
Two fixtures cover most needs:
\begin{itemize}[leftmargin=*]
  \item \texttt{client} returns a \texttt{TestClient} bound to the app;
  \item \texttt{auth\_headers} registers a user, logs in, and returns the
        \texttt{Authorization} header with the resulting bearer token.
\end{itemize}

\section{Sample Test}
\begin{lstlisting}[style=base, language=Python]
def test_create_student_with_auth(client, auth_headers):
    resp = client.post(
        "/students/",
        json={"name": "Ada Lovelace", "department": "Math"},
        headers=auth_headers,
    )
    assert resp.status_code == 201
    body = resp.json()
    assert body["name"] == "Ada Lovelace"
\end{lstlisting}

\section{What to Test}
A pragmatic checklist: each endpoint's happy path, each authentication
boundary, validation on input edge cases, and at least one negative case
per error branch (\texttt{401}, \texttt{404}, \texttt{422}).

% =========================================================================
\chapter{GitHub Workflow}
\label{ch:github}

The repository ships with a single \texttt{main} branch. New work is
done on a feature branch:
\begin{lstlisting}[style=base]
git checkout -b feature/<short-name>
# ... edits ...
git add -A && git commit -m "feat: <imperative summary>"
git push -u origin feature/<short-name>
\end{lstlisting}
A pull request describes \emph{why} the change exists, not just what it
does. Reviews focus on correctness, naming, and the test that pins the
behavior in place.

\section{Continuous Integration (Optional)}
A minimal GitHub Actions workflow can run the test suite on every push.
The job installs the requirements, exports a SQLite \texttt{DATABASE\_URL},
and invokes \texttt{pytest -q}.

\section{Releases}
Tag releases with semantic versions (\texttt{v1.0.0}). The container image
can be built and pushed from the same workflow on tags only.

% =========================================================================
\chapter{Conclusion and Next Steps}
\label{ch:next}

The repository is intentionally small, but it covers the moving parts of
a production-grade Python web service. Reasonable next steps include:
\begin{itemize}[leftmargin=*]
  \item replacing \texttt{create\_all} with Alembic migrations;
  \item adding role-based authorization to the JWT payload;
  \item adding rate limiting and structured logging;
  \item shipping a CI pipeline that builds, tests, and publishes images;
  \item exposing more domain resources (courses, enrollments, grades).
\end{itemize}

The companion repository will continue to evolve, and so should this
book. The shortest distance between an idea and a running endpoint is
still a pull request.

\backmatter

\end{document}
